%% toc: P1 3 3.7,3.8
\section{Proof for the Ultrafilter Lemma at its first statement}
\label{sec:annex-p01-w02b-p01-005-ultrafilter-lemma-entry}

Throughout, (F1)--(F4) and (D) are the conditions labelled in
\cref{sec:p01-w02-p01-009-assumptions}, and the ambient theory is
\(\mathsf{ZF}\) unless a step names \(\mathsf{AC}\).

\begin{annexitem}{annex:p01-w02b-p01-005-ultrafilter-lemma}{thm:p01-w02b-p01-005-ultrafilter-lemma}{Ultrafilter Lemma, with its standing}
\begin{proof}
\emph{Identity of the statements.} The hypotheses of part~1 are the set
\(I\) and conditions (F1)--(F4) on \(\mathcal F\subseteq\mathcal P(I)\); its
conclusion is a family \(\mathcal U\subseteq\mathcal P(I)\) with
\(\mathcal F\subseteq\mathcal U\), (F1)--(F4) and (D). Part~1 of
Theorem~\ref{thm:p01-w02-p01-009-ultrafilter-lemma} has exactly these
hypotheses and this conclusion, under the same assumption
\(\mathsf{ZF}+\mathsf{AC}\). Part~2 here and part~2 there are the same
\(\mathsf{ZF}\)-sentence, and part~3 here and part~4 there are the same
assertion about provability under the same hypothesis. The arguments below
are written out for this statement.

\emph{Part 1.} Assume \(\mathsf{ZF}+\mathsf{AC}\), and let \(I\) and
\(\mathcal F\) satisfy (F1)--(F4). By Lemma~\ref{thm:p01-w02-p01-009-terms},
part~1, whose proof uses no choice principle, \(\mathcal F\) is a proper
Boolean filter of the Boolean algebra \(\mathcal P(I)\) with meet \(\cap\),
join \(\cup\), complement \(X\mapsto I\setminus X\), bottom \(\varnothing\)
and top \(I\). Theorem~\ref{thm:p01-bpi-ac-implication}, proved in
Annex~\ref{annex:p01-bpi-ac-implication} from \(\mathsf{AC}\) through Zorn's
lemma (Theorem~\ref{thm:p06-ac-to-zorn}), gives in
\(\mathsf{ZF}+\mathsf{AC}\) an ultrafilter of \(\mathcal P(I)\) containing
\(\mathcal F\), that is, a proper Boolean filter
\(\mathcal U\supseteq\mathcal F\) maximal under inclusion among proper
Boolean filters of \(\mathcal P(I)\). By
Lemma~\ref{thm:p01-w02-p01-009-terms}, part~2, such a \(\mathcal U\)
satisfies (F1)--(F4) and (D). Since \(\mathcal F\subseteq\mathcal U\), this
\(\mathcal U\) is the family that part~1 asserts to exist.

\emph{Part 2.} The derivation of part~1 is a derivation in the first-order
theory \(\mathsf{ZF}\) with the single additional axiom \(\mathsf{AC}\), and
\(I\) and \(\mathcal F\) are arbitrary in it, so the universal closure of the
principle over \(I\) and \(\mathcal F\) is derivable in
\(\mathsf{ZF}+\mathsf{AC}\). Since \(\mathsf{AC}\) is a sentence, the
deduction theorem for first-order logic turns this into a derivation in
\(\mathsf{ZF}\) of the implication from \(\mathsf{AC}\) to that universal
closure.

\emph{Part 3.} Write \(\mathsf{UL}\) for the universal closure of the
principle of part~1, and assume that \(\mathsf{ZF}\) is consistent. Suppose,
for a contradiction, that \(\mathsf{ZF}\) proves
\(\mathsf{UL}\rightarrow\mathsf{AC}\).

By Theorem~\ref{thm:p01-w02-p01-009-ultrafilter-lemma}, part~3, proved in
\(\mathsf{ZF}\) in Annex~\ref{annex:p01-w02-p01-009-lemma} from
Lemma~\ref{thm:p01-w02-p01-009-terms} and Theorem~\ref{thm:p01-bpi-equivalence},
\(\mathsf{ZF}\) proves \(\mathsf{BPI}\rightarrow\mathsf{UL}\). Chaining the
two implications inside \(\mathsf{ZF}\), \(\mathsf{ZF}\) proves
\(\mathsf{BPI}\rightarrow\mathsf{AC}\).

Let \(T=\mathsf{ZF}+\mathsf{BPI}+\neg\mathsf{AC}\). Every theorem of
\(\mathsf{ZF}\) is a theorem of \(T\), so \(T\) proves
\(\mathsf{BPI}\rightarrow\mathsf{AC}\); \(T\) proves \(\mathsf{BPI}\) as an
axiom, hence \(\mathsf{AC}\) by modus ponens; and \(T\) proves
\(\neg\mathsf{AC}\) as an axiom. So \(T\) proves a sentence and its negation,
and \(T\) is inconsistent.

The theorem of Halpern and L\'evy, recorded as
Theorem~\ref{thm:p01-bpi-choice-separation} \cite{HalpernLevy1971}, states
that \(\operatorname{Con}(\mathsf{ZF})\) implies \(\operatorname{Con}(T)\).
Since \(\mathsf{ZF}\) was assumed consistent, \(T\) is consistent, which
contradicts the previous paragraph. Hence \(\mathsf{ZF}\) does not prove
\(\mathsf{UL}\rightarrow\mathsf{AC}\).

The only premise of this argument whose proof is not given in this handbook
is the Halpern--L\'evy relative-consistency theorem; its model construction is
cited, not reproduced. Every other step is a first-order step about
provability or an appeal to a result proved in this handbook.

\emph{The concluding sentence.} Part~2 gives the \(\mathsf{ZF}\)-provable
implication \(\mathsf{AC}\rightarrow\mathsf{UL}\), and part~3 gives, if
\(\mathsf{ZF}\) is consistent, the \(\mathsf{ZF}\)-unprovability of
\(\mathsf{UL}\rightarrow\mathsf{AC}\). Their conjunction is the stated sense
of ``strictly weaker over \(\mathsf{ZF}\)''. In particular, under that
hypothesis \(\mathsf{ZF}\) does not prove
\(\mathsf{UL}\leftrightarrow\mathsf{AC}\), so the principle is not a second
name for the Axiom of Choice. That \(\mathsf{UL}\) is a sentence of the
ambient set theory, quantifying over sets \(I\) and families
\(\mathcal F\subseteq\mathcal P(I)\), and not a formula of an object theory,
is immediate from its form.
\end{proof}
\end{annexitem}
